% ============================================================
% 基礎演習Ⅱ　産学連携プロジェクト　企業への質問リスト
%   テーマ「2050年に向けた新店舗設計」（カーディーラー）について、
%   企業担当者との質疑応答で学生が使う質問の例（32問）と、
%   自分の質問・回答のメモ欄。
%   デザインは各回ワークシート（../worksheets/preamble.tex）と共通。
%   ビルド：make（Makefile を参照）
% ============================================================
% 共通のプリアンブル（worksheets.tex と session.tex から読み込む）
\documentclass[paper=a4paper,fontsize=10pt,line_length=48zw,number_of_lines=44,baselineskip=1.6zw]{jlreq}

\usepackage{xcolor}
\usepackage[most]{tcolorbox}
\usepackage{tabularray}
\usepackage{enumitem}
\usepackage{needspace}
\usepackage{fancyhdr}
\usepackage[hidelinks]{hyperref}

% ------------------------------------------------------------
% 色（サイトの assets/css/style.css に合わせる）
% ------------------------------------------------------------
\definecolor{ink}{HTML}{1D2330}
\definecolor{inktwo}{HTML}{4A5261}
\definecolor{inkthree}{HTML}{757D8A}
\definecolor{rulegray}{HTML}{B9BEB4}
\definecolor{papertwo}{HTML}{ECEEE8}
\definecolor{pone}{HTML}{2747C9}  \definecolor{ponetint}{HTML}{DCE2F9}
\definecolor{ptwo}{HTML}{D68A17}  \definecolor{ptwotint}{HTML}{F6E4C3}
\definecolor{pthree}{HTML}{6E7684}\definecolor{pthreetint}{HTML}{E0E3E4}
\definecolor{bad}{HTML}{B3261E}   \definecolor{badtint}{HTML}{F9DEDA}
\colorlet{phase}{pone}
\colorlet{phasetint}{ponetint}

% ------------------------------------------------------------
% ページスタイル
% ------------------------------------------------------------
\newcommand*\sessionmark{}
\pagestyle{fancy}
\fancyhf{}
\renewcommand{\headrulewidth}{0pt}
\fancyfoot[L]{\sffamily\scriptsize\color{inkthree}基礎演習Ⅱ　2026年度後期　ワークシート}
\fancyfoot[R]{\sffamily\scriptsize\color{inkthree}\sessionmark\quad\thepage}
\fancypagestyle{cover}{\fancyhf{}}

\setlength{\parindent}{0pt}
\setlist{nosep,leftmargin=1.5em}

% ------------------------------------------------------------
% 氏名欄（記入例・模範解答版では answers-preamble.tex で中身を入れる）
\newcommand*\wsid{}
\newcommand*\wsname{}
\newcommand*\wsgroup{}
\newcommand\fieldline[2]{\makebox[#1][l]{\rlap{\hspace{1mm}#2}\hrulefill}}

% 各回の見出し
%   \session{回}{フェーズ 1|2|3}{日付}{題目}{タグ}{グループ欄 yes|no}
% ------------------------------------------------------------
\newcommand\session[6]{%
  \clearpage
  \ifcase#2\or
    \colorlet{phase}{pone}\colorlet{phasetint}{ponetint}\def\phasename{探究の基礎}\or
    \colorlet{phase}{ptwo}\colorlet{phasetint}{ptwotint}\def\phasename{実践プロジェクト}\or
    \colorlet{phase}{pthree}\colorlet{phasetint}{pthreetint}\def\phasename{キャリア講習}\fi
  \renewcommand*\sessionmark{第#1回　#4}%
  \phantomsection\addcontentsline{toc}{section}{\texorpdfstring{第#1回　#4\hfill{\small #3}}{第#1回　#4}}%
  \begin{tcolorbox}[enhanced,colback=phase,colframe=phase,arc=2mm,boxrule=0pt,
      left=4mm,right=4mm,top=2.5mm,bottom=2.5mm,fontupper=\color{white}\sffamily]
    {\small 第#1回　#3　\textbar　\phasename\if\relax\detokenize{#5}\relax\else　\textbar　#5\fi}\par
    \vspace{0.6mm}{\LARGE\bfseries #4}
  \end{tcolorbox}
  \vspace{-1mm}
  {\sffamily\small\color{inktwo}
  学籍番号\fieldline{30mm}{\wsid}\quad 氏名\fieldline{42mm}{\wsname}%
  \ifx\relax#6\relax\else\quad グループ\fieldline{28mm}{\wsgroup}\fi}\par
  \vspace{2mm}%
}

% ねらい
\newcommand\goals[1]{%
  \begin{tcolorbox}[enhanced,colback=phasetint,colframe=phasetint,arc=1.5mm,boxrule=0pt,
      left=3mm,right=3mm,top=1.5mm,bottom=1.5mm,fontupper=\small]
    {\sffamily\bfseries\color{phase}この回のねらい}\par
    \begin{itemize}[label=\textcolor{phase}{・},leftmargin=1.2em]#1\end{itemize}
  \end{tcolorbox}}

% ワークの見出し　\work[必要な高さ]{ラベル}{題}
\newcommand\work[3][36mm]{%
  \par\needspace{#1}\vspace{2.5mm}%
  {\sffamily\bfseries\color{phase}■ #2}\enspace{\sffamily\bfseries #3}\par\vspace{1mm}}

% 指示文
\newcommand\instr[1]{{\small\color{inktwo}#1}\par\vspace{1mm}}

% 記入欄　\answerbox[欄内の小見出し]{高さ}
\newcommand\answerbox[2][]{%
  \begin{tcolorbox}[enhanced,colback=white,colframe=rulegray,boxrule=0.5pt,arc=1.2mm,
      height=#2,left=2mm,right=2mm,top=1.2mm,bottom=1mm,
      fontupper=\sffamily\scriptsize\color{inkthree},before skip=1mm,after skip=1.5mm]#1\end{tcolorbox}}

% 横並びの記入欄　\twoboxes{見出し1}{見出し2}{高さ}
\newcommand\twoboxes[3]{%
  \par\noindent
  \begin{minipage}[t]{0.49\linewidth}\answerbox[#1]{#3}\end{minipage}\hfill
  \begin{minipage}[t]{0.49\linewidth}\answerbox[#2]{#3}\end{minipage}\par}

% 方眼の記入欄（スケッチ用）
\newcommand\gridbox[2][]{%
  \begin{tcolorbox}[enhanced,colback=white,colframe=rulegray,boxrule=0.5pt,arc=1.2mm,
      height=#2,left=2mm,top=1.2mm,fontupper=\sffamily\scriptsize\color{inkthree},
      underlay={\begin{tcbclipinterior}\draw[step=5mm,black!12,very thin]
        (interior.south west) grid (interior.north east);\end{tcbclipinterior}}]#1\end{tcolorbox}}

% 罫線　\writelines{行数}
\newcommand\writelines[1]{%
  \par\foreach\n in {1,...,#1}{\nointerlineskip\noindent
    {\color{rulegray}\vbox to 7.5mm{\vfil\hrule width\linewidth}}\par}\vspace{1.5mm}}

% 1行の記入欄　\fillin{見出し}
\newcommand\fillin[1]{%
  \par\noindent{\sffamily\small #1}\enspace{\color{rulegray}\leaders\hrule\hfill}\par\vspace{2.5mm}}

% プロンプト例
\newtcolorbox{promptbox}[1][]{enhanced,colback=papertwo,colframe=papertwo,boxrule=0pt,
  arc=1mm,leftrule=1.2mm,colbacktitle=papertwo,
  borderline west={1.2mm}{0pt}{phase},left=3mm,right=2mm,top=1mm,bottom=1mm,
  fontupper=\sffamily\footnotesize\color{inktwo},before skip=1mm,after skip=2mm,
  title={\sffamily\scriptsize\bfseries\color{phase}プロンプト例　#1},
  attach title to upper={\par\vspace{0.5mm}}}

% 注意
\newtcolorbox{caution}[1]{enhanced,colback=badtint,colframe=badtint,boxrule=0pt,arc=1mm,
  left=3mm,right=3mm,top=1mm,bottom=1mm,fontupper=\footnotesize,
  before upper={{\sffamily\bfseries\color{bad}#1}\par}}

% チェック項目
\newenvironment{checks}{\begin{itemize}[label={□},leftmargin=1.5em,itemsep=0.8mm]\small}{\end{itemize}}

% 次回までに
\newenvironment{nexttime}{%
  \par\needspace{24mm}\vspace{2.5mm}%
  \begin{tcolorbox}[enhanced,colback=white,colframe=phase,boxrule=0.6pt,arc=1.5mm,
      left=3mm,right=3mm,top=1.5mm,bottom=1.5mm,fontupper=\small]
    {\sffamily\bfseries\color{phase}次回までに}\par\vspace{0.5mm}
    \begin{itemize}[label={□},leftmargin=1.5em,itemsep=0.6mm]}%
  {\end{itemize}\end{tcolorbox}}

% AI活用記録（その回の分）
\newcommand\aimemo{%
  \par\needspace{50mm}\vspace{2.5mm}%
  {\sffamily\bfseries\color{phase}■ AI活用記録（この回の分）}\par\vspace{1mm}
  \begin{wstab}{colspec={Q[l,wd=30mm,m] X[l,m]},row{1-Z}={ht=9mm}}
    使用した場面 & \\
    主なプロンプト & \\
    出力をどう扱ったか\newline{\scriptsize 採用・修正・不採用と理由} & \\
    検証したこと\newline{\scriptsize 何で確かめたか} & \\
  \end{wstab}\par}

% 表の共通スタイル
\NewTblrEnviron{wstab}
\SetTblrInner[wstab]{hlines={0.4pt,rulegray},vlines={0.4pt,rulegray},
  cells={font=\small},rowsep=2pt,colsep=4pt}
\NewTblrEnviron{wstabh}
\SetTblrInner[wstabh]{hlines={0.4pt,rulegray},vlines={0.4pt,rulegray},
  cells={font=\small},rowsep=2pt,colsep=4pt,
  row{1}={bg=phasetint,font=\sffamily\footnotesize\bfseries,ht=0pt}}

% ペルソナカード（第6回）
\newcommand\personacard[1]{%
\begin{wstab}{colspec={Q[l,wd=26mm,m] X[l,m]},rows={ht=7.5mm}}
  \SetCell[c=2]{bg=phasetint,font=\sffamily\footnotesize\bfseries}ペルソナ#1　名前： & \\
  年齢・職業 & \\ 家族構成 & \\ 居住地の特徴 & \\ 来店の目的 & \\ 重視する点 & \\
  不安や面倒 & \\ 休日の過ごし方 & \\
  \SetRow{ht=14mm}本人の言葉\newline\scriptsize 店舗に求めること & \\
\end{wstab}}

% 見出しセル（左列を色付きにする用）
\newcommand\hc[1]{\SetCell{bg=phasetint,font=\sffamily\footnotesize\bfseries}#1}


\colorlet{phase}{ptwo}
\colorlet{phasetint}{ptwotint}
\fancyfoot[L]{\sffamily\scriptsize\color{inkthree}基礎演習Ⅱ　2026年度後期　企業への質問リスト}
\renewcommand*\sessionmark{2050年に向けた新店舗設計}

% カテゴリの見出し　\qcat{記号}{題}{説明}
\newcommand\qcat[3]{%
  \par\needspace{36mm}\vspace{2mm}%
  {\sffamily\bfseries\color{phase}■ #1}\enspace{\sffamily\bfseries #2}\par
  \vspace{0.3mm}{\footnotesize\color{inktwo}#3}\par\vspace{0.8mm}}

% 質問　\q{質問文}{ねらい}
\newcounter{qno}
\newcommand\q[2]{%
  \stepcounter{qno}%
  \par\needspace{12mm}%
  {\leftskip=12.5mm\noindent
    \llap{\makebox[12.5mm][l]{\sffamily\bfseries\color{phase}□\,Q\theqno}}%
    {\small #1}\par
    \vspace{0.2mm}%
    {\footnotesize\color{inktwo}\textsf{ねらい}\enspace #2}\par}%
  \vspace{1.2mm}}

\begin{document}

% ------------------------------------------------------------
% 見出し
% ------------------------------------------------------------
\begin{tcolorbox}[enhanced,colback=phase,colframe=phase,arc=2mm,boxrule=0pt,
    left=4mm,right=4mm,top=2.5mm,bottom=2.5mm,fontupper=\color{white}\sffamily]
  {\small 産学連携プロジェクト　\textbar　企業担当者との質疑応答}\par
  \vspace{0.6mm}{\LARGE\bfseries 企業への質問リスト}\par
  \vspace{0.3mm}{\normalsize テーマ：2050年に向けた新店舗設計}
\end{tcolorbox}
\vspace{-1mm}
{\sffamily\small\color{inktwo}
学籍番号\fieldline{30mm}{}\quad 氏名\fieldline{42mm}{}\quad グループ\fieldline{28mm}{}}\par
\vspace{2mm}

\begin{tcolorbox}[enhanced,colback=phasetint,colframe=phasetint,arc=1.5mm,boxrule=0pt,
    left=3mm,right=3mm,top=1.5mm,bottom=1.5mm,fontupper=\small]
  {\sffamily\bfseries\color{phase}このリストの使い方}\par
  \begin{itemize}[label=\textcolor{phase}{・},leftmargin=1.2em]
    \item 質問の例を、A〜Hの8つの観点で32問あげています。そのまま使っても、自分たちの提案に合わせて言い換えてもかまいません。
    \item 事前にグループで読み、\textbf{聞きたい質問の□に印}をつけておきます。グループ内で質問が重ならないよう、誰がどれを聞くかを決めておきましょう。
    \item 「ねらい」は、その質問の答えが提案のどこに役立つかを示しています。答えを聞いたら、最後のページのメモ欄に記録します。
    \item 質問の時間は限られています。自分たちの提案に直結する質問を優先し、資料や調べればわかることは聞かないようにします。
  \end{itemize}
\end{tcolorbox}

\begin{tcolorbox}[enhanced,colback=white,colframe=rulegray,boxrule=0.5pt,arc=1.5mm,
    left=3mm,right=3mm,top=1.5mm,bottom=1.5mm,fontupper=\small]
  {\sffamily\bfseries\color{phase}よい質問にするコツ}\par
  \begin{itemize}[label=\textcolor{phase}{・},leftmargin=1.2em]
    \item \textbf{最初に名乗り、質問の背景を一言添える}：「○○グループの△△です。私たちは□□という店舗を考えていて、……」
    \item \textbf{一度に聞くのは1つ}：複数を聞きたいときは「2点あります」と最初に数を言う。
    \item \textbf{「はい・いいえ」で終わらない聞き方にする}：「〜はありますか」より「〜はどのように考えていますか」。
    \item \textbf{答えを深める}：「具体的には、どのような場面ですか」「それはなぜですか」「例えば、どんなお客様ですか」。
    \item \textbf{最後にお礼を言う}：「ありがとうございます。提案に生かします」。
  \end{itemize}
\end{tcolorbox}

\begin{caution}{企業の情報の取り扱い}
  質疑応答で聞いた内容のうち、非公開のもの（数字、社内の事情、検討中の計画など）は、AIに入力しない、SNSや外部に書かない、グループ外に共有しないことを守ってください。判断に迷う場合は担当教員に確認します。
\end{caution}

% ------------------------------------------------------------
% 質問
% ------------------------------------------------------------
\qcat{A}{いまの店舗と課題を知る}{2050年を考える前に、出発点となる現在の姿をつかむ質問です。}

\q{現在の店舗で、お客様が来店してから帰るまでの流れのうち、最も時間がかかっている場面や、不満が出やすい場面はどこですか。}
  {設計で解決すべき具体的な課題を見つける。}
\q{来店されるお客様は、購入の相談、点検・修理、納車など、どのような目的で来られる方が多いですか。ここ10年で目的の内訳はどう変わりましたか。}
  {店舗のどの機能にどれだけ面積を割くかを考える手がかりにする。}
\q{店舗を設計・改装するときに、自動車メーカーのブランド基準など、ディーラーとして自由に変えられない決まりはありますか。}
  {提案が守るべき前提条件を確かめる。}
\q{いまの店舗で、お客様から特に評価されている点と、社内で「ここは変えたい」と話題になっている点を教えてください。}
  {残すべき強みと、改善すべき弱みを分けて把握する。}

\qcat{B}{2050年のお客様とクルマ}{25年後に、誰が、どのようにクルマと関わっているかを想像するための質問です。}

\q{2050年には、クルマを「所有する」人と、サブスクリプションやカーシェアで「利用する」人の割合はどう変わると見ていますか。それによって店舗の役割はどう変わりますか。}
  {「売る場所」から何に変わるのか、店舗の存在意義を考える。}
\q{自動運転が普及した場合、お客様がディーラーに求めるものは何になるとお考えですか。}
  {運転しない時代の顧客ニーズをペルソナに反映する。}
\q{2050年の主なお客様は、いまの10代・20代です。私たちの世代に、どのような体験や関わり方を期待されていますか。}
  {学生ならではの視点を、企業の期待と結びつける。}
\q{高齢化が進むなかで、運転免許を返納された方やそのご家族とは、どのような関係を続けていきたいとお考えですか。}
  {クルマを運転しない人も含めた顧客像を考える。}
\q{オンラインで車を買う人が増えても、それでもわざわざ店舗に来る理由として、何が残ると思いますか。}
  {リアルの店舗でしか提供できない価値を見きわめる。}

\qcat{C}{店舗の役割と機能}{2050年の店舗に「何があるべきか」を考えるための質問です。}

\q{2050年の店舗は、車を売る以外にどのような役割を担うべきだとお考えですか（地域の交流拠点、エネルギーの拠点、移動サービスの窓口など）。}
  {コンセプトの方向性について、企業の考えを聞く。}
\q{電気自動車や水素自動車が主流になると、整備・点検の内容や、サービス工場の広さ・設備はどう変わりますか。}
  {ショールームと工場の面積配分や配置を考える根拠にする。}
\q{展示車を置かない、または台数を大きく減らした店舗は考えられますか。その場合、車の魅力をどのように伝えますか。}
  {大胆な空間構成がどこまで受け入れられるかを確かめる。}
\q{車を買う予定のない人にも来てもらえる店舗は、ビジネスとして成り立つと思いますか。成り立たせるには何が必要ですか。}
  {「人が集まる場」のアイデアを、事業として説明できるようにする。}

\qcat{D}{空間とデザイン}{立地、動線、居心地、だれもが使いやすいことなど、設計そのものに関わる質問です。}

\q{店舗の立地は、今後も郊外の幹線道路沿いが中心になるとお考えですか。駅前、商業施設の中、小さな店舗を分散させる形なども検討されていますか。}
  {提案する店舗の立地と規模の前提を決める。}
\q{お客様が店舗で長い時間を過ごすのは、商談、待ち時間など、どの場面ですか。その時間を快適にするために試した工夫と、その結果を教えてください。}
  {待合・商談スペースの設計に、実際の経験を生かす。}
\q{子ども連れ、高齢の方、車いすを使う方、外国籍の方など、さまざまなお客様が使いやすい店舗にするうえで、いまの課題は何ですか。}
  {ユニバーサルデザインの観点を提案に取り入れる。}
\q{外観や内装に「ディーラーらしさ」を残すべきでしょうか。それとも、まったく別の業態のように見えてもよいでしょうか。}
  {デザインの自由度と、ブランドとして守るべき印象を確かめる。}

\qcat{E}{テクノロジーとデジタル}{オンラインやAIと、リアルの店舗との役割分担を考えるための質問です。}

\q{VR・ARによる試乗体験や、オンライン上の店舗での商談を、どのように評価されていますか。リアルの店舗との役割分担はどうなると思いますか。}
  {デジタルとリアルを組み合わせた来店体験を設計する。}
\q{AIによる接客や、スタッフのいない時間帯がある店舗は、将来ありうるとお考えですか。人が対応すべきこととして何が残りますか。}
  {人の接客が生きる場面に、空間を集中させる。}
\q{来店前にお客様がオンラインで得た情報と、店舗での体験をつなげるために、いまどのような工夫をされていますか（差し支えない範囲で）。}
  {来店前から来店後までの一連の流れ（顧客体験）を描く。}
\q{通信でつながる車（コネクテッドカー）が普及すると、点検のお知らせやアフターサービスはどう変わりますか。それは店舗の設計にも影響しますか。}
  {アフターサービスの拠点としての店舗の姿を考える。}

\qcat{F}{環境と地域社会}{カーボンニュートラルや防災など、社会の中での店舗の役割を考えるための質問です。}

\q{2050年のカーボンニュートラルに向けて、店舗の建物そのもの（太陽光発電、蓄電池、断熱、木造など）で取り組みたいことはありますか。}
  {環境に配慮した建物の提案に、企業の方針を反映する。}
\q{災害のときに、店舗や電気自動車を地域の非常用電源や避難場所として役立てる構想はありますか。}
  {地域に必要とされる店舗という価値を加える。}
\q{店舗がある地域の特徴（人口の変化、公共交通、気候など）は、店舗づくりにどの程度反映されていますか。}
  {地域に合わせた店舗にするための調べ方を知る。}
\q{地域の住民、学校、自治体などと連携した取り組みで、これまでにうまくいったものがあれば教えてください。}
  {実績のある取り組みを手がかりに、地域との関わりを提案する。}

\qcat{G}{店舗で働く人}{お客様だけでなく、そこで働く人にとってよい店舗かどうかを考える質問です。}

\q{2050年の店舗で働く人には、どのような役割やスキルが求められるとお考えですか。}
  {スタッフの役割に合わせた空間（相談、体験、整備など）を考える。}
\q{従業員の方の働きやすさ（休憩スペース、動線、バックヤードなど）について、いまの店舗の課題は何ですか。提案に含めてもよいですか。}
  {お客様から見えない部分まで含めて設計する。}
\q{人手不足が続くなかで、少ない人数でも質の高い接客を続けるために、店舗の設計で工夫できることは何だと思いますか。}
  {省力化と接客の質を両立する動線やレイアウトを考える。}

\qcat{H}{提案に向けて}{提案の方向性と評価の観点を確かめるための質問です。最後に聞くとよいでしょう。}

\q{私たちの提案に、特に期待していることは何ですか。実現できるかどうかと、大胆な発想のどちらを重視して評価されますか。}
  {提案の力点（現実性と独創性のバランス）を決める。}
\q{予算、敷地の広さ、法律による規制などの条件は、どの程度意識して提案すればよいですか。想定すべき店舗の規模や立地があれば教えてください。}
  {提案の前提条件をそろえ、絵に描いた餅にならないようにする。}
\q{これまでに社内などで出た新店舗のアイデアのうち、魅力的だったけれど実現しなかったものはありますか。実現しなかった理由も教えてください。}
  {実現の壁になりやすい点を知り、提案で先回りして答える。}
\q{提案をまとめる途中で、追加の質問や相談をすることはできますか。できる場合、どのような方法がよいですか。}
  {質疑応答のあとも、企業とやりとりできる道をつくる。}

% ------------------------------------------------------------
% メモ欄
% ------------------------------------------------------------
\clearpage
{\sffamily\bfseries\Large\color{phase}質問と回答のメモ}\par\vspace{1mm}
\instr{上のリストから選んだ質問や、自分たちで考えた質問を書き、質疑応答の場で聞いた回答をメモします。非公開の情報の扱いに注意してください。}

\begin{wstabh}{colspec={Q[c,wd=12mm,m] X[2,l,m] X[3,l,m]},row{2-Z}={ht=22mm}}
  番号 & 質問（だれが聞くか） & 回答のメモ \\
  & & \\
  & & \\
  & & \\
  & & \\
  & & \\
  & & \\
  & & \\
  & & \\
\end{wstabh}\par

\answerbox[質疑応答を聞いて、提案に生かしたいこと・グループで話し合いたいこと]{30mm}

\end{document}
